% Section 7 (writer: steep-a; edited by p2-body). ex:split:noisyparity and rem:split:siprograms are in app-split.tex.
% Source: research/tracks/steeper/notes-final.md §3.1-§3.4, Remark 4.3.
% Renamed from the source: samples N -> n, input length bound n_N -> d_n, Ville deficiency m* -> delta*, Hellinger H^2 -> \Hel.
\section{The model--randomness split}\label{sec:split}

A second question of the author's notes is whether a steeper prior prefers a simple sampler to programs in which the sampled bits are filled in; at $\lambda=1$, Solomonoff induction has no such preference. This section makes both halves precise. At $\lambda=1$ the sampler and the programs that fill in its sampled bits get similar posterior, up to explicit factors: $2^{k+\delta^*}$ in one direction and $2^{2\log(t_n+1)+\kappa}$ in the other, where $2\log(t_n+1)$ grows at most triply logarithmically in $n$ and in the sampled randomness (not at all if the sampler's label probabilities are bounded below) (\cref{thm:split:lambda1,prop:split:si}). At fixed $\lambda>1$ with private randomness samplers beat fill-ins (not necessarily $\nu$ itself; \cref{rem:split:reading}), in every countable product class (\cref{thm:split:private}). With shared randomness, or in SI, the universal mixture wins instead and no split is made (\cref{thm:split:shared}). Proofs are in \cref{app:split}.

\subsection{Setting}\label{sec:split:setting}

\begin{definition}[Samplers and fill-ins]\src{steeper \S3.1}\label{def:split:setting}
A \emph{sampler} is a computable conditional measure $\nu$ ($\nu(y\mid x)$ computable uniformly in $(x,y)$, $\sum_y\nu(y\mid x)=1$) with a program $\pi_\nu$ of length $k$ in the class considered. Distinct inputs $x_1,x_2,\dots$ are chosen by any rule, possibly adaptive, randomised or non-computable, and then $y_i\sim\nu(\cdot\mid x_i)$; $\Prob$ is the joint law and $D_n$ the first $n$ pairs. The \emph{sampled randomness} is $-\log\nu(D_n)$. \emph{Fill-ins} are deterministic programs consistent with $D_n$; in a product class their total weight is $\Wdet(D):=\sum\{2^{-\ell(\pi)}:\pi\text{ deterministic},\ \nu_\pi\supseteq D\}$.
\end{definition}

The first lemma bounds how much any program, or the whole mixture, can gain over the sampler on its own data.

\begin{lemma}[Ville deficiency]\status{proved; computed; Ville's inequality known}\src{steeper Lemma 3.1; check c4(A); referee q3(B)}\label{lem:split:ville}
Let $\xi$ be a shared hypothesis with $\xi(\emptyset)\le1$, for example the mixture $\xi_{\mathcal C}$ of a countable product class $\mathcal C$ with $\sum_{\pi\in\mathcal C}2^{-\ell(\pi)}\le1$ (for $\Hcl{\times}$ this is $\Mprod$). Let $\delta^*:=\sup_n\log\big(\xi(D_n)/\nu(D_n)\big)$. Then $\Prob[\delta^*\ge j]\le2^{-j}$ for all $j\ge0$. Consequently, for $\xi=\xi_{\mathcal C}$, every $\pi\in\mathcal C$ and every $n$,
\[
\nu_\pi(D_n)\le2^{\ell(\pi)+\delta^*}\nu(D_n).
\]
Adaptive inputs are allowed.
\end{lemma}

\emph{Proof idea.} $\xi(D_n)/\nu(D_n)$ is a nonnegative supermartingale with initial value at most $1$, by the one-point inequality; apply Ville's maximal inequality \citep{ville1939etude} (\cref{app:split:ville}).

Fill-ins, on the other hand, are cheap: a deterministic program can store a finite piece of each input's private random stream. Adapting the tape-length bound of \citet{claudeSFI2026fundamentals} % P1: prop:versions:tape
to product classes and to every computable sampler gives, for every data set $D$ with $n$ inputs of length at most $d$ and $\beta_D:=\min_{(x,y)\in D}\nu(y\mid x)>0$, the bound $\Wdet(D)\ge2^{-k-2\log(t_D+1)-\kappa}\nu(D)$, where $2\log(t_D+1)$ is doubly logarithmic in $d$ and triply logarithmic in $n$ and in the sampled randomness $-\log\nu(D)$ (\cref{lem:split:stream} in \cref{app:split:stream}). This needs a computable \emph{measure}; for dyadic samplers a hashing argument gives a weaker bound (\cref{lem:split:hash}).

\subsection{$\lambda=1$: the sampler and its fill-ins get similar posterior}\label{sec:split:lambda1}

For a run of the data process put $d_n:=\max_{i\le n}|x_i|$, $\beta_n:=\min_{i\le n}\nu(y_i\mid x_i)$, $i_n:=\lceil\log(6n/\beta_n)\rceil$ and $t_n:=\lceil\log(2d_n+2+2\log(i_n+1))\rceil$. Let $\Pi_1(\det\mid D_n)$ be the posterior mass of all deterministic programs.

\begin{theorem}[SFI at $\lambda=1$]\status{proved}\src{steeper Thm 3}\label{thm:split:lambda1}
In the setting of \cref{def:split:setting} with the class $\Hcl{\times}$, for every sampler $\nu$ with a program of length $k$, simultaneously for all $n$:
\begin{enumerate}[label=(\roman*)]
\item $\Pi_1(\pi_\nu\mid D_n)\ge2^{-k-\delta^*}$;
\item $\Pi_1(\det\mid D_n)=\Wdet(D_n)/\xi_\times(D_n)\ge2^{-k-2\log(t_n+1)-\kappa-\delta^*}$;
\item the posterior odds of all fill-ins against the sampler program lie in $\big[2^{-2\log(t_n+1)-\kappa},\,2^{k+\delta^*}\big]$, with a lower end that depends on the sampled randomness only through $t_n$;
\item for a sampler whose values are multiples of $2^{-r}$, \cref{lem:split:hash} gives the alternative lower ends $2^{-k-2\log(n+1)-2\log(d_n+1)-2\log(r+1)-\kappa-\delta^*}$ in (ii) and $2^{-2\log(n+1)-2\log(d_n+1)-2\log(r+1)-\kappa}$ in (iii).
\end{enumerate}
\end{theorem}

\begin{proof}
$\Pi_1(\pi_\nu\mid D_n)=2^{-k}\nu(D_n)/\xi_\times(D_n)$ and $\Pi_1(\det\mid D_n)=\Wdet(D_n)/\xi_\times(D_n)$. Bound $\xi_\times(D_n)\le2^{\delta^*}\nu(D_n)$ by \cref{lem:split:ville}, $\Wdet(D_n)$ from below by \cref{lem:split:stream} (or \cref{lem:split:hash} for (iv)) and from above by $\xi_\times(D_n)$.
\end{proof}

The lower ends of (ii) and (iii) hold for every data set, not only for samples from $\nu$, and cover non-dyadic samplers: one that gives every label probability at least $1/3$, such as the Bern$(1/3)$ threshold hypotheses in Sam Eisenstat's coin example (personal communication; see also \citet{claudeSFI2026fundamentals}), has $i_n\le\lceil\log(18n)\rceil$. The hashing bound of (iv) is weaker at realistic sizes (32 to 143 bits; \cref{app:split:hash}).

In SI, the total weight of the deterministic hypotheses consistent with the data $z_{1:n}$ is at least $2^{-k-K(n)-\kappa}\nu(z_{1:n})$, and in SI's representation by deterministic programs every hypothesis is a fill-in (\cref{prop:split:si} in \cref{app:split:si}). So at $\lambda=1$ the sampler and its fill-ins get similar posterior in SFI and in SI, up to factors $2^k$, $2^{2\log(t_n+1)+\kappa}$ (SFI) or $2^{O(k+K(n))}$ (SI), and the Ville deficiency, with the dependence on the sampled randomness only through $t_n$.

\subsection{$\lambda>1$ with private randomness: samplers beat fill-ins}\label{sec:split:private}

Call a program $\pi$ \emph{$\nu$-equivalent on the run} if $\sum_n\Hel\big(\nu(\cdot\mid x_n),\nu_\pi(\cdot\mid x_n)\big)<\infty$, where $\Hel(P,Q):=1-\sum_y\sqrt{P(y)Q(y)}\in[0,1]$, also for a sub-probability $Q$.

\begin{theorem}[Private randomness, $\lambda>1$]\status{proved; computed; Kakutani-type argument known}\src{steeper Thm 4; check c4(C); referee r4, q3(C)}\label{thm:split:private}
Let $\mathcal C$ be a countable class of product programs with $\sum_{\pi\in\mathcal C}2^{-\ell(\pi)}\le1$ that contains $\pi_\nu$ (for example $\Hcl{\times}$, or a time-bounded or circuit subclass). Fix $\lambda>1$ and define $\delta^*$ from $\xi_{\mathcal C}$. Simultaneously for all $n$:
\begin{enumerate}[label=(\roman*)]
\item \emph{MAP.} Every $\lambda$-MAP $\hat\pi_n$ has $\ell(\hat\pi_n)\le(\lambda k+\delta^*)/(\lambda-1)=k+(k+\delta^*)/(\lambda-1)$ and $-\log\nu(D_n)-\ell(\hat\pi_n)-\delta^*\le-\log\nu_{\hat\pi_n}(D_n)\le-\log\nu(D_n)+\lambda k$.
\item \emph{Length tail.} $\Pi_\lambda(\{\pi:\ell(\pi)>L\}\mid D_n)\le2^{\lambda k+\delta^*-(\lambda-1)L}$ for every $L$.
\item \emph{Fill-ins.} Every deterministic program consistent with $D_n$ has $\ell\ge-\log\nu(D_n)-\delta^*$; $\Pi_\lambda(\det\mid D_n)\le2^{\lambda k+\lambda\delta^*-(\lambda-1)(-\log\nu(D_n))}$; and for the tempered posterior with $\eta=1/\lambda$, $\Pi^\eta(\det\mid D_n)\le2^{k+\delta^*-(1-\eta)(-\log\nu(D_n))}$.
\item \emph{Identification.} Suppose the input sequence is infinite. Then, almost surely, from some $n$ on every $\lambda$-MAP is $\nu$-equivalent on the run, and $\Pi_\lambda(\nu\text{-equivalent on the run}\mid D_n)\to1$.
\end{enumerate}
\end{theorem}

\emph{Proof idea.} (i)--(iii) compare with the sampler program, $J_\lambda(\pi_\nu)=\lambda k-\log\nu(D_n)$, and bound every program's likelihood by \cref{lem:split:ville}. For (iv), the likelihood ratio $\Lambda_n:=\nu_\pi(D_n)/\nu(D_n)$ of a program that is not $\nu$-equivalent on the run tends to $0$ almost surely, by the martingale argument behind Kakutani's dichotomy \citep{kakutani1948equivalence}, which holds for adaptive inputs (\cref{app:split:private}).

\begin{remark}[What the theorem says, and what it does not]\status{remark}\src{steeper \S3.3}\label{rem:split:reading}
This is the split asked for in \cref{sec:intro}, for every fixed $\lambda>1$, in a precise and limited sense. The selected model is a sampler, not a fill-in, but not necessarily $\nu$: it may carry up to $(k+\delta^*)/(\lambda-1)$ bits beyond the sampler, for example memorised noise. In particular no $\lambda$-MAP is a fill-in once $-\log\nu(D_n)>\lambda(k+\delta^*)/(\lambda-1)$. The fill-ins' posterior is exponentially small in the sampled randomness, but positive. ``$\nu$-equivalent on the run'' is not equality of predictive laws and says nothing about unseen inputs. On a finite domain such as $\{0,1\}^d$, (iv) is vacuous.
\end{remark}

\subsection{$\lambda>1$ with shared randomness: no split}\label{sec:split:shared}

\begin{theorem}[SI, or SFI with shared randomness, $\lambda>1$]\status{proved; computed}\src{steeper Thm 5; check c4(C); referee r4, q4(A)}\label{thm:split:shared}
Assume (Dom) and $\xi_1\le2^{\kappa_1}U$, let $\lambda>1$, and let $\pi_\nu$ be a program of length $k$ for a sampler in the class. For every observation $o$ with $U(o)>0$:
\begin{enumerate}[label=(\roman*)]
\item if $(\lambda-1)k>\lambda\ell_U+\kappa_0$, then $J_\lambda(\pi_U;o)<J_\lambda(\pi_\nu;o)$: the sampler program is never a $\lambda$-MAP, and every $\lambda$-MAP has length at most $\ell^*_\lambda<k$;
\item for every set $A$ of programs, $\Pi_\lambda(A\mid o)\le2^{-(\lambda-1)\min_{\pi\in A}\ell(\pi)+\lambda\ell_U+\kappa_1}$. In particular the power posterior puts at most $2^{-(\lambda-1)L+\lambda\ell_U+\kappa_1}$ on all programs of length at least $L$.
\end{enumerate}
For SI with adaptive or non-computable inputs the sampler is not an SI hypothesis, since it does not model the inputs; use $\Hcl{\mathrm{chr}}$, where it is a chronological semimeasure and (Dom) holds with $U=\MH$. With i.i.d.\ computable inputs, $\Hcl{\mathrm{SI}}$ with $U=M$ works.
\end{theorem}

\emph{Proof idea.} (i) compares $\pi_U$ with $\pi_\nu$ using (Dom); (ii) bounds the numerator by $2^{-(\lambda-1)\min_A\ell}\xi_1(o)$ and the normaliser by $2^{-\lambda\ell_U}U(o)$ (\cref{app:split:shared}).

Compare \cref{thm:split:private}. In a shared class the steeper posterior lives, up to mass $2^{-s}$, on programs of length at most $(\lambda\ell_U+\kappa_1+s)/(\lambda-1)$, \emph{whatever the data}; in the product class it concentrates on samplers $\nu$-equivalent on the run, whose length is unbounded over $\nu$. For fixed $\lambda$, both steeper predictors dominate $\nu$ ($\xi^\times_\lambda\ge2^{-\lambda k}\nu$; $\xi^F_\lambda\ge2^{-\lambda\ell_U-\kappa_0-k}\nu$), so by the argument of the merging theorem of \citet{claudeSFI2026fundamentals}% P1: thm:learn:merging
\ their summed expected KL divergences to $\nu$'s predictions are at most $\lambda k$ and $\lambda\ell_U+\kappa_0+k$. The difference is \emph{what the selected hypothesis is}: a model plus private noise, or a short universal learner that contains no model of $\nu$.

\Cref{app:split:shared} illustrates the contrast on a small noisy parity (\cref{ex:split:noisyparity}) and shows that in SI's representation by deterministic programs no form of steepness separates a sampler from its fill-ins (\cref{rem:split:siprograms}).
